\chapter{Debugging}

\hayashi{} includes a built-in Debug Adapter Protocol (DAP) server. With the VS~Code extension you can set breakpoints, step through a script, and inspect variables.

\section{Starting a debug session}

Make sure the \hay{} binary is in \texttt{PATH}. In VS~Code open a \texttt{.hay} file, go to the \textbf{Run and Debug} panel, choose the \textbf{"Debug Hayashi script"} configuration, and start debugging.

Behind the scenes VS~Code runs:

\begin{lstlisting}[style=inline, language=sh]
hay dap script.hay
\end{lstlisting}

The DAP server supports breakpoints, step over/in/out, continue, scopes (Locals / Globals), and variable inspection.

\section{Inspecting model results}

When execution pauses, model objects expand in the Variables panel with a concise summary and structured children.

\begin{lstlisting}[style=inline]
result: OLS(k=2, n=10000), R2=1.0000
  coefficients    DataFrame(2 rows, 7 cols)
  fit             Dict(13 entries)
  residuals       Series(residuals: 10000 values)
  fitted_values   Series(fitted_values: 10000 values)
  params          Series(params: 2 values)
  std_errors      Series(std_errors: 2 values)
  test_values     Series(test_values: 2 values)
  p_values        Series(p_values: 2 values)
  conf_lower      Series(conf_lower: 2 values)
  conf_upper      Series(conf_upper: 2 values)
\end{lstlisting}

\subsection{Children}

\begin{description}
  \item[\texttt{coefficients}] Coefficient table as a \texttt{DataFrame}: \texttt{variable}, \texttt{coef}, \texttt{std\_err}, \texttt{t} or \texttt{z}, \texttt{p\_value}, \texttt{conf\_low}, \texttt{conf\_high}.
  \item[\texttt{fit}] Model-specific fit statistics as a \texttt{Dict} (R$^2$, pseudo R$^2$, log-likelihood, AIC/BIC, \texttt{n\_obs}, \texttt{df\_resid}, etc.).
  \item[\texttt{params}, \texttt{std\_errors}, \texttt{test\_values}, \texttt{p\_values}] Raw coefficient vectors as \texttt{Series}.
  \item[\texttt{conf\_lower} / \texttt{conf\_upper}] Confidence interval bounds, when available.
  \item[\texttt{residuals} / \texttt{fitted\_values}] OLS only.
\end{description}

\section{Supported estimators}

The structured expansion works for the following result types:

\begin{itemize}
  \item Cross-section: OLS, IV/2SLS, quantile regression, GMM, GLSAR
  \item Panel: fixed effects, random effects, PCSE, PanelGLS, FE-2SLS, Arellano-Bond, system GMM
  \item Binary: logit and probit
  \item Count: Poisson, negative binomial, zero-inflated Poisson/negative binomial
  \item Limited: Tobit, ordered logit/probit, multinomial logit
  \item Other GLM: GLM, RLM, beta regression, mixed-effects / HLM
  \item Systems: SUR and three-stage least squares (each equation is a child node)
\end{itemize}

\begin{nota}
  Time-series models (ARIMA, GARCH, VAR, etc.) currently display as plain text in the debugger without structured children.
\end{nota}

\section{Command-line debugging}

You can start the DAP server manually for testing or integration with other editors:

\begin{lstlisting}[style=inline, language=sh]
hay dap script.hay
\end{lstlisting}

It reads DAP messages from \texttt{stdin} and writes responses to \texttt{stdout}. The VS~Code extension handles the protocol automatically.
